%% kapa: exam=1 round=37 year=2026 period=1 form=A subject=economics subject_ko=경제학원론 seq=34 number=74
%% source: Q-Net 감정평가사 기출문제 (cst00309) artlSeq=5258924 file=raw/1008/5258924/2026년 제37회 감정평가사 1차 1교시.pdf
%% conversion: kordoc 4.24.51 -> tools/gen1.py (auto); review=proofread-sonnet; flags=none
\begin{question}{74}
거시경제 안정화정책의 시차에 관한 설명으로 옳은 것을 모두 고른 것은?
\begin{framedtext}
ㄱ. 금융정책은 재정정책에 비해 내부시차가 더 짧다. \par
ㄴ. 소득세와 같은 자동안정화 장치는 내부시차를 줄여준다. \par
ㄷ. 외부시차는 경제충격 발생 시점과 충격에 대응한 정책시행 간의 시차를 뜻한다. \par
\end{framedtext}
\begin{choices}
  \item ㄱ
  \item ㄱ, ㄴ
  \item ㄱ, ㄷ
  \item ㄴ, ㄷ
  \item ㄱ, ㄴ, ㄷ
\end{choices}
\end{question}
\answer{2}
